\documentclass[10pt,a4paper]{amsart}
\usepackage[margin=19mm]{geometry}
\usepackage{amsmath,amssymb,mathtools}
\usepackage[scr=rsfs,cal=euler]{mathalfa}
\usepackage{xfrac}
\usepackage[unicode,hidelinks,bookmarks=false]{hyperref}
\newcommand{\ie}{i.e.\ }
\newcommand{\cf}{cf.\ }
\newcommand{\resp}{respectively\ }
\newcommand{\Subsection}{\subsection}
\providecommand{\nobreakdash}{\nobreak\hbox{-}}
\setlength{\emergencystretch}{2em}
\multlinegap0pt
\usepackage{bm}
\usepackage{bbm}
\usepackage{dsfont}
\usepackage{xspace}
\usepackage{cancel}
\usepackage{mathtools}
\usepackage{amsthm}
\makeatletter
\@ifundefined{theorem}{\newtheorem{theorem}{Theorem}}{}
\@ifundefined{lemma}{\newtheorem{lemma}[theorem]{Lemma}}{}
\@ifundefined{proposition}{\newtheorem{proposition}[theorem]{Proposition}}{}
\makeatother
\newcommand{\Om}{\Omega}
\newcommand{\R}{\mathbb R}
\newcommand{\Sn}{\mathbb S^{n-1}}
\newcommand{\tr}{\operatorname{tr}}
\title[An inverse source problem for the Monge--Ampere equation fro]{An inverse source problem for the Monge--Ampere equation from large boundary data}
\author{C\u{a}t\u{a}lin I. C\^arstea, Tuhin Ghosh}
\date{Source: arXiv:2606.07064v1 (CC BY 4.0)}
\begin{document}
\maketitle

\noindent\emph{Source and licence.} An inverse source problem for the Monge--Ampere equation from large boundary data. Version arXiv:2606.07064v1 (CC BY 4.0), distributed under the Creative Commons Attribution 4.0 International licence, as recorded in the arXiv licence metadata for this exact version and shown on the abstract page https://arxiv.org/abs/2606.07064v1. Full rights notice: licenses/arxiv-2606.07064-CC-BY-4.0.txt.\par\smallskip

\noindent\textit{Editorial scope.} Counted unit: the Proposition whose source label is \texttt{prop:boundary-normal-from-C0} in the article (source0.tex lines 1189--1211, source label \texttt{prop:boundary-normal-from-C0}), delivered together with the article's own complete original proof of it (source0.tex lines 1213--1375, one \texttt{proof} environment of 163 lines). No mathematics is written, repaired or re-derived: statement and proof are the article's own text line for line. Required context retained uncounted: lem:non-glancing-regularity (lines 482--513); lem:det-first-order (lines 753--793); lem:schur-convexity-criterion (lines 884--915); prop:comparison-barriers (lines 1033--1068); eq:p-x0-normal-zero (lines 1160--1164); eq:b-normal-step8 (lines 1175--1180). Every definition, display and label of those lines is retained unchanged. Omitted from this package and not redistributed: 1--481, 514--752, 794--883, 916--1032, 1069--1159, 1165--1174, 1181--1188, 1212--1212, 1376--1696. Nothing inside a retained excerpt is omitted or altered: every retained body is delivered exactly as the article's lines are written, so no omission receipt of any kind accompanies this package. Citations: the retained excerpts carry no citation command at all, so no bibliography entry of the article is transcribed and no reference is redistributed. Reference adaptation: none was needed and none was made; every reference inside the delivered unit resolves inside it, so no unresolved reference or citation is produced. Printed slips kept literal: no printed word, symbol, hypothesis or number of the counted statement or of its proof was altered. Layout and class: the article's own class is \texttt{article} with packages including geometry, amsmath, amssymb, amsthm, mathtools, mathrsfs, enumitem, hyperref; the wrapper is the lane's pinned \texttt{amsart} editorial wrapper at 10pt on A4, which restates the theorem-like environments the retained text needs and adds the article's own macro definitions that the retained text needs (\texttt{\textbackslash Om}, \texttt{\textbackslash R}, \texttt{\textbackslash Sn}, \texttt{\textbackslash tr}), with extra wrapper packages (bm, bbm, dsfont, xspace, cancel, mathtools). The delivered numbering is the unit's own. Assets: no retained span contains a figure environment, a graphics-inclusion command, a PDF-page inclusion, a table or a TikZ picture; no image file is copied into this package, the package declares no asset dependency and no image file is redistributed with it. Licence: version arXiv:2606.07064v1 of this article is distributed under the Creative Commons Attribution 4.0 International licence, as recorded in the arXiv licence metadata for this exact version and shown on the abstract page https://arxiv.org/abs/2606.07064v1. A full rights notice is delivered at licenses/arxiv-2606.07064-CC-BY-4.0.txt. Characters: every retained excerpt is pure ASCII, so the delivered engine, XeLaTeX with its default Latin Modern text font, reports no missing character.
\begin{lemma}\label{lem:non-glancing-regularity}
Let \(m\ge0\) and \(0<\kappa<1\).  There is a constant
\(C=C(\Omega,m,\kappa)\), independent of \(\omega\), such that the following hold.

First,
\begin{equation}\label{eq:s-pm-Cm-estimates}
    \|s_+\|_{C^m(\Omega_{\omega,\kappa})}
    +
    \|s_-\|_{C^m(\Omega_{\omega,\kappa})}
    +
    \|\ell_\omega^{-1}\|_{C^m(\Omega_{\omega,\kappa})}
    \le C .
\end{equation}
Second, if \(f\in C^m(\overline\Om)\), then
\begin{equation}\label{eq:w-Cm-nonglancing-estimate}
    \|w_\omega^f\|_{C^m(Q_{\omega,\kappa})}
    \le
    C\|f\|_{C^m(\overline\Om)} .
\end{equation}
Here the \(C^m\)-norm on \(Q_{\omega,\kappa}\) is taken in the coordinates \((y,s)\in\omega^\perp\times\R\).  In particular, for every such non-glancing region,
\begin{equation}\label{eq:w-C2-nonglancing-estimate}
    \|D^2 w_\omega^f\|_{L^\infty(Q_{\omega,\kappa})}
    \le
    C\|f\|_{C^2(\overline\Om)} .
\end{equation}
Moreover, after increasing the constant if necessary,
\begin{equation}\label{eq:normal-derivative-nonglancing-estimate}
    \|\partial_\nu w_\omega^f\|_{C^m(\Gamma_{\omega,\kappa}^+\cup\Gamma_{\omega,\kappa}^-)}
    \le
    C\|f\|_{C^{m+1}(\overline\Om)} .
\end{equation}
\end{lemma}

\begin{lemma}\label{lem:det-first-order}
Let \(M\ge1\).  There exists a constant \(C=C(n,M)\) such that the following holds.  Suppose
\[
    |a|+|z|+\|B\|\le M,
    \qquad
    t^{-n}\|B\|\le \frac12.
\]
Then
\begin{equation}\label{eq:det-first-order}
\begin{aligned}
    \det D^2(t\phi_\omega+t^{1-n}w)
    &=
    a
    +
    t^{-n}\left(a\tr B-|z|^2\right)
    +
    R_t,
\end{aligned}
\end{equation}
where
\begin{equation}\label{eq:det-remainder-bound}
    |R_t|\le C t^{-2n}.
\end{equation}
Equivalently,
\begin{equation}\label{eq:det-first-order-w}
\begin{aligned}
    \det D^2(t\phi_\omega+t^{1-n}w)
    &=
    \partial_s^2 w \\
    &\quad
    +
    t^{-n}\left[
        (\partial_s^2w)\tr(D_y^2w)
        -
        |D_y\partial_s w|^2
    \right]
    +O(t^{-2n}).
\end{aligned}
\end{equation}
The constant in the \(O(t^{-2n})\) term depends only on \(n\) and on the stated bound for the second derivatives of \(w\).
\end{lemma}

\begin{lemma}\label{lem:schur-convexity-criterion}
Let \(c_0>0\), \(M<\infty\), and \(L<\infty\).  Suppose that on \(Q\)
\begin{equation}\label{eq:convexity-hypotheses-general}
    \partial_s^2 w\ge c_0,
    \qquad
    \|D^2w\|+\|D^2\eta\|\le M .
\end{equation}
Then there is \(T=T(n,c_0,M,L)\) such that, for every \(t\ge T\) and every \(|\lambda|\le L\),
\[
    D^2V_{t,\lambda}>0
    \quad\text{on }Q.
\]
More precisely, writing
\[
    W_{t,\lambda}=w+\lambda t^{-n}\eta,
\]
one has
\begin{equation}\label{eq:convexity-transverse-block}
    I+t^{-n}D_y^2W_{t,\lambda}
    \ge \frac12 I
\end{equation}
and
\begin{equation}\label{eq:convexity-schur-lower}
    \partial_s^2W_{t,\lambda}
    -t^{-n}
    \left\langle
        (I+t^{-n}D_y^2W_{t,\lambda})^{-1}D_y\partial_sW_{t,\lambda},
        D_y\partial_sW_{t,\lambda}
    \right\rangle
    \ge \frac{c_0}{2} .
\end{equation}
\end{lemma}

\begin{proposition}\label{prop:comparison-barriers}
Let \(f\in C^\infty(\overline\Omega)\) satisfy \(f\ge c_f>0\), and fix \(\omega\in\Sn\).  There exist constants \(C_*>0\) and \(T<\infty\), depending on \(f\), \(\omega\), and \(\Omega\), such that, for every \(t\ge T\), the functions
\begin{equation}\label{eq:lower-upper-barriers}
    V_t^{\rm low}
    =U_t+C_*t^{1-2n}b_\omega,
    \qquad
    V_t^{\rm up}
    =U_t-C_*t^{1-2n}b_\omega
\end{equation}
are strictly convex in \(\Omega\), agree with \(t\phi_\omega\) on \(\partial\Omega\), and satisfy
\begin{equation}\label{eq:barrier-order-around-U}
    V_t^{\rm low}\le U_t\le V_t^{\rm up}
    \quad\text{in }\Omega,
\end{equation}
\begin{equation}\label{eq:barrier-det-inequalities}
    \det D^2V_t^{\rm low}\ge f,
    \qquad
    \det D^2V_t^{\rm up}\le f
    \quad\text{in }\Omega .
\end{equation}
Consequently,
\begin{equation}\label{eq:global-C0-asymptotic}
    U_t+C_*t^{1-2n}b_\omega
    \le
    u_{t,\omega}^f
    \le
    U_t-C_*t^{1-2n}b_\omega
    \quad\text{in }\Omega,
\end{equation}
and hence
\begin{equation}\label{eq:global-C0-asymptotic-bound}
    \|u_{t,\omega}^f-U_t\|_{L^\infty(\Omega)}
    \le
    \frac{C_*}{8}\operatorname{diam}(\Omega)^2 t^{1-2n}.
\end{equation}
\end{proposition}

\begin{equation}\label{eq:p-x0-normal-zero}
    p_{x_0}(x_0)=0,
    \qquad
    \partial_\nu p_{x_0}(x_0)=0.
\end{equation}

\begin{equation}\label{eq:b-normal-step8}
    \partial_\nu b_\omega(x_\pm(y))
    =
    \frac{\ell_\omega(y)}{2|\omega\cdot\nu(x_\pm(y))|}
    >0.
\end{equation}

\begin{proposition}\label{prop:boundary-normal-from-C0}
Let \(f\in C^\infty(\overline\Omega)\) satisfy \(f\ge c_f>0\).  Fix \(\omega\in\Sn\), and let \(K\Subset\Gamma_\omega\).  Then there are constants \(C_K<\infty\) and \(T_K<\infty\) such that, for all \(t\ge T_K\),
\begin{equation}\label{eq:boundary-normal-rate-from-C0}
    \sup_{x\in K}
    \left|
        \partial_\nu u_{t,\omega}^f(x)
        -\partial_\nu U_t(x)
    \right|
    \le
    C_K t^{1-n-n/2}.
\end{equation}
In particular,
\begin{equation}\label{eq:boundary-normal-asymptotic-proved}
    \partial_\nu u_{t,\omega}^f
    =
    t\partial_\nu\phi_\omega
    +
    t^{1-n}\partial_\nu w_\omega^f
    +
    o(t^{1-n})
\end{equation}
locally uniformly on \(\Gamma_\omega\).
\end{proposition}

\begin{proof}
We prove the estimate pointwise with constants uniform for \(x_0\in K\).  Write
\[
    K_+=K\cap\Gamma_\omega^+,
    \qquad
    K_-=K\cap\Gamma_\omega^- .
\]
Since \(K\Subset\Gamma_\omega\), there are \(\kappa>0\) and an open set \(Y\Subset\Omega_\omega\) such that, with
\[
    Y_+=P_{\omega^\perp}(K_+),
    \qquad
    Y_-=P_{\omega^\perp}(K_-),
\]
one has \(Y_+\cup Y_-\subset Y\) and
\[
    \omega\cdot\nu(x_+(y))\ge\kappa,
    \qquad
    -\omega\cdot\nu(x_-(y))\ge\kappa
    \qquad \text{for all }y\in Y .
\]
Lemma~\ref{lem:non-glancing-regularity} gives uniform bounds for the endpoint functions and for \(\partial_\nu b_\omega\) on this non-glancing set.  Since \(\overline Y\subset\Omega_\omega\) and \(\ell_\omega\) is continuous and positive on \(\Omega_\omega\), there is \(\ell_0>0\) such that
\[
    \ell_\omega(y)=s_+(y)-s_-(y)\ge \ell_0
    \qquad \text{for all }y\in Y .
\]
Choose \(\rho>0\) so small that \(\{y\in\omega^\perp:|y-y_0|<\rho\}\subset Y\) for every \(y_0\in Y_+\cup Y_-\), and then choose \(0<\delta<\ell_0/4\).  Reducing \(\rho\) and \(\delta\), if necessary, we may also assume that the following collars remain inside this non-glancing tube.  If \(x_0\in K_+\), set \(y_0=P_{\omega^\perp}x_0\) and
\begin{equation}\label{eq:plus-boundary-collar}
    D_{x_0}^+
    =
    \{y+s\omega: |y-y_0|<\rho,\ s_+(y)-\delta<s<s_+(y)\}.
\end{equation}
If \(x_0\in K_-\), set
\begin{equation}\label{eq:minus-boundary-collar}
    D_{x_0}^-
    =
    \{y+s\omega: |y-y_0|<\rho,\ s_-(y)<s<s_-(y)+\delta\}.
\end{equation}
We write \(D_{x_0}\) for either collar.  The Monge--Ampere comparison principle for convex solutions is valid on bounded domains, so the corners of these collars cause no difficulty; one may also justify the same step by applying the comparison principle on a smooth exhaustion and passing to the limit.

Let \(\varepsilon_t\) and \(\eta_t\) be as above.  Define
\begin{equation}\label{eq:boundary-local-barriers}
    L_{t,x_0}=U_t-\eta_tG_{x_0},
    \qquad
    W_{t,x_0}=U_t+\eta_tG_{x_0}.
\end{equation}
Since \(G_{x_0}\ge0\), these barriers are ordered around \(U_t\).  On the physical boundary part of \(D_{x_0}\), one has \(w_\omega^f=0\) and \(b_\omega=0\), so
\[
    L_{t,x_0}=t\phi_\omega-\eta_t p_{x_0}
    \le
    t\phi_\omega
    =u_{t,\omega}^f
    \le
    t\phi_\omega+\eta_t p_{x_0}
    =W_{t,x_0}.
\]
At the distinguished point \(x_0\), both inequalities are equalities.

On the artificial boundary of \(D_{x_0}\), the function \(G_{x_0}\) is bounded below by a positive constant independent of \(x_0\).  On the tangential face \(|y-y_0|=\rho\), one has \(G_{x_0}\ge p_{x_0}=\rho^2\).  On the inner face, for either sign,
\[
    -b_\omega=\frac12\delta(\ell_\omega(y)-\delta)
    \ge \frac12\delta(\ell_0-\delta).
\]
Thus there exists \(c_0>0\) such that
\begin{equation}\label{eq:G-positive-artificial-boundary}
    G_{x_0}\ge c_0
    \quad\text{on }\partial D_{x_0}\setminus\partial\Omega,
\end{equation}
for every \(x_0\in K\).

By Proposition~\ref{prop:comparison-barriers}, there is \(M<\infty\) such that
\begin{equation}\label{eq:global-C0-input-for-boundary}
    \|u_{t,\omega}^f-U_t\|_{L^\infty(\Omega)}\le M t^{1-2n}
\end{equation}
for all sufficiently large \(t\).  Since
\[
    \frac{\eta_t}{t^{1-2n}}=t^{n/2}\to\infty,
\]
\eqref{eq:G-positive-artificial-boundary} and \eqref{eq:global-C0-input-for-boundary} imply, for all sufficiently large \(t\),
\[
    L_{t,x_0}\le u_{t,\omega}^f\le W_{t,x_0}
    \quad\text{on }\partial D_{x_0}.
\]

We next check the determinant inequalities.  The profiles in \eqref{eq:boundary-local-barriers} are
\[
    L_{t,x_0}
    =
    t\phi_\omega+t^{1-n}(w_\omega^f-\varepsilon_tG_{x_0}),
    \qquad
    W_{t,x_0}
    =
    t\phi_\omega+t^{1-n}(w_\omega^f+\varepsilon_tG_{x_0}).
\]
On the chosen collars, \(G_{x_0}\) has uniformly bounded second derivatives and \(\partial_s^2G_{x_0}=-1\).  Lemma~\ref{lem:det-first-order} gives, uniformly for \(x\in D_{x_0}\) and \(x_0\in K\),
\begin{equation}\label{eq:boundary-barrier-det-expansion}
\begin{aligned}
    \det D^2L_{t,x_0}&=f+\varepsilon_t+O(t^{-n}),\\
    \det D^2W_{t,x_0}&=f-\varepsilon_t+O(t^{-n}).
\end{aligned}
\end{equation}
Since \(\varepsilon_t=t^{-n/2}\) and \(t^{-n}=o(\varepsilon_t)\), after increasing the lower threshold for \(t\),
\begin{equation}\label{eq:boundary-barrier-det-inequalities}
    \det D^2L_{t,x_0}\ge f,
    \qquad
    \det D^2W_{t,x_0}\le f
    \quad\text{in }D_{x_0}.
\end{equation}
The same uniform \(C^2\)-bounds and the Schur-complement convexity argument of Lemma~\ref{lem:schur-convexity-criterion} show that both \(L_{t,x_0}\) and \(W_{t,x_0}\) are strictly convex for all sufficiently large \(t\).  Indeed, their longitudinal profile second derivatives are \(f+\varepsilon_t\) and \(f-\varepsilon_t\), both bounded below by \(c_f/2\) for large \(t\), while the transverse block remains a small relative perturbation of \(tI_{n-1}\).

The comparison principle, with the orientation recalled above, now gives
\begin{equation}\label{eq:boundary-barrier-comparison}
    L_{t,x_0}\le u_{t,\omega}^f\le W_{t,x_0}
    \quad\text{in }D_{x_0}.
\end{equation}
Indeed, the lower barrier has determinant at least \(f=\det D^2u_{t,\omega}^f\), so it lies below \(u_{t,\omega}^f\); the upper barrier has determinant at most \(f\), so \(u_{t,\omega}^f\) lies below it.

At \(x_0\), the three functions in \eqref{eq:boundary-barrier-comparison} have the same boundary value.  Since \(u_{t,\omega}^f-L_{t,x_0}\ge0\) and \(W_{t,x_0}-u_{t,\omega}^f\ge0\) in the one-sided collar, their outward normal derivatives at \(x_0\) are nonpositive.  Hence
\begin{equation}\label{eq:normal-derivative-squeeze}
    \partial_\nu W_{t,x_0}(x_0)
    \le
    \partial_\nu u_{t,\omega}^f(x_0)
    \le
    \partial_\nu L_{t,x_0}(x_0).
\end{equation}
Using \eqref{eq:p-x0-normal-zero} and \(G_{x_0}=-b_\omega+p_{x_0}\), we get
\[
    \partial_\nu G_{x_0}(x_0)=-\partial_\nu b_\omega(x_0).
\]
Therefore
\[
    \partial_\nu L_{t,x_0}(x_0)
    =
    \partial_\nu U_t(x_0)
    +\eta_t\partial_\nu b_\omega(x_0),
    \qquad
    \partial_\nu W_{t,x_0}(x_0)
    =
    \partial_\nu U_t(x_0)
    -\eta_t\partial_\nu b_\omega(x_0).
\]
By \eqref{eq:b-normal-step8}, \(\partial_\nu b_\omega(x_0)>0\), and the non-glancing estimates give
\[
    \sup_K |\partial_\nu b_\omega|\le C_K .
\]
Thus \eqref{eq:normal-derivative-squeeze} implies
\[
    \left|
        \partial_\nu u_{t,\omega}^f(x_0)
        -\partial_\nu U_t(x_0)
    \right|
    \le C_K\eta_t
    =C_Kt^{1-n-n/2},
\]
uniformly for \(x_0\in K\).  Finally,
\[
    \partial_\nu U_t
    =
    t\partial_\nu\phi_\omega
    +
    t^{1-n}\partial_\nu w_\omega^f,
\]
and \(t^{1-n-n/2}=o(t^{1-n})\).  This proves the proposition.
\end{proof}

\end{document}
